% ============================================================================
% SECTION 2.2 : LE CADRE CLO — FORMALISATION DES MESURES
% Source : partie2/clo.tex (version condensée, niveau de détail maximal)
% Preuves complètes intégrées dans le corps (pas de renvoi aux annexes)
% ============================================================================

\section{LE CADRE CLO : FORMALISATION DES MESURES}
\label{sec:clo}
\begin{sloppypar}
Le chapitre~\ref{chap:contexte} a expos\'e le paradoxe de l'authenticit\'e
invisible~\cite{Goodman15,Casey16,Grimmelmann15}, selon lequel les protocoles qui maximisent la protection cryptographique
induisent une opacit\'e institutionnelle s\'ev\`ere, tandis que les protocoles
transparents sacrifient la confidentialit\'e~\cite{Bonneau16,Citron14}. La section~\ref{sec:etat_art}
a fourni les outils n\'ecessaires, th\'eorie de l'information de
Shannon~\cite{Shannon48,Shannon49,Wyner75}, in\'egalit\'es de Pinsker et
Fano~\cite{Fano61,Pinsker1964}, th\'eorie de l'information quantique
(von Neumann, Holevo~\cite{Holevo1973,NielsenChuang2010}), primitives d'obscurcissement (commitments,
zk-SNARKs~\cite{Groth16,BCTV14,Gennaro13}), trilemmes d'impossibilit\'e (CAP~\cite{Brewer00},
FLP~\cite{FLP85}), standards post-quantiques NIST~\cite{NISTPQC,Kyber2022,Dilithium2022}, sans encore les articuler 
en un cadre unifi\'e.

Nous d\'efinissons ici ce cadre. Le \textbf{cadre CLO (Cryptographic
Legal Opposability)}~\cite{MINKA2026CLO} constitue la premi\`ere formalisation
math\'ematique rigoureuse de la cryptographie institutionnelle. L'
Objectif~\ref{obj:obj1_unification}, comment \'etablir un cadre unifi\'e
formalisant confidentialit\'e, fiabilit\'e et opposabilit\'e via des mesures
issues de la th\'eorie de l'information permettant analyses quantitatives et comparaisons
objectives~? re\c{c}oit une r\'eponse en cinq \'etapes, qui portent successivement sur le
mod\`ele de protocole, les d\'efinitions des m\'etriques, les propri\'et\'es
fondamentales, le paradigme de conception et une construction prototype.

Notre apport conceptuel central est l'opposabilit\'e $\Opp$. La
litt\'erature cryptographique classique traite
l'interpr\'etabilit\'e institutionnelle comme contrainte externe~\cite{GMR85,Goldreich01,GoldwasserMicali84} ou lacune
\'educative des utilisateurs~\cite{Whitten99,Adams1999,Herley09}. Le cadre CLO en fait une dimension \`a
optimiser au m\^eme titre que $\Conf$ et $\Rel$.

Cette section correspond \`a l'\'etape \textbf{F (Formaliser)} de la progression \textbf{F}$\to$B$\to$E$\to$S$\to$C$\to$X et apporte la contribution \textbf{C1} de cette th\`ese.
\end{sloppypar}

\begin{boiteArbitre}
L'arbitre dispose d\'esormais d'une m\'etrique pour ce qu'il
ressentait. L'opposabilit\'e $\mathrm{Opp} = I(V;\Sigma|C)/H(V|C)$ mesure
exactement la fraction d'incertitude institutionnelle que la preuve
cryptographique r\'esout. Selon les estimations illustratives du
cadre, Helios lui laisse $50\%$ d'incertitude r\'esiduelle~; ECDSA lui en
laisse $10\%$. Il peut maintenant comparer et se demander ce qui, dans le
protocole ou dans l'institution, limite les compromis accessibles.
\end{boiteArbitre}

% -----------------------------------------------------------------------
\subsection{Mod\`ele de protocole CLO}
\label{subsec:modele_protocole_clo}
% -----------------------------------------------------------------------
\begin{sloppypar}

\subsubsection{D\'efinition formelle}
\label{subsubsec:definition_protocole}

Un protocole d'evidence cryptographique op\'erant en contexte institutionnel
doit satisfaire deux niveaux de v\'erification fondamentalement distincts,
\`a savoir la correction math\'ematique (v\'erification cryptographique) et l'extraction
de signification s\'emantique (interpr\'etation institutionnelle)~\cite{ESIGN00,eIDAS14,Kerr04}. Cette
dualit\'e, absente des d\'efinitions cryptographiques
classiques~\cite{GMR85,Goldreich01,GoldwasserMicali84} motive le mod\`ele formel qui suit.

\begin{definition}[Protocole CLO]\label{def:protocole_clo}
Un \textbf{protocole CLO} $\mathcal{P}$ consiste en quatre espaces et
quatre algorithmes param\'etr\'es par~$\lambda$.

\medskip\noindent\textit{Espaces.}
Le protocole manipule l'espace de messages $\mathcal{M}_\lambda$ (typiquement $\{0,1\}^\lambda$),
l'espace de preuves cryptographiques $\Sigma_\lambda$,
l'espace de contextes institutionnels $\mathcal{C}_\lambda$
et l'espace de verdicts $\mathcal{V}_\lambda$ (d\'ecisions institutionnelles).

\medskip\noindent\textit{Algorithmes.}
\begin{itemize}
\item $\mathsf{Setup}(1^\lambda) \to (pk, sk)$ : g\'en\'eration de cl\'es publique/secr\`ete
\item $\mathsf{Prove}(sk, m, c) \to \sigma$ : g\'en\'eration de preuve pour message $m$ et contexte $c$
\item $\mathsf{Verify}(pk, \sigma, c) \to \{0,1\}$ : v\'erification cryptographique de la preuve
\item $\mathsf{Interpret}(\sigma, c) \to v$ : \textbf{interpr\'etation institutionnelle} produisant verdict $v \in \mathcal{V}$
\end{itemize}

\medskip\noindent o\`u $c$ repr\'esente la distribution du contexte
institutionnel capturant l'environnement d\'ecisionnel (pr\'ec\'edents l\'egaux,
expertise disponible, proc\'edures \'etablies).
\end{definition}

\begin{remark}[Double couche de v\'erification]\label{rem:double_couche}
$\mathsf{Verify}$ r\'epond \`a la question \emph{\guillemotleft~La preuve est-elle math\'ematiquement
valide~?\guillemotright} $\mathsf{Interpret}$ r\'epond \`a la question \emph{\guillemotleft~Que signifie cette
preuve pour ma d\'ecision~?\guillemotright} La tension centrale de la cryptographie
institutionnelle \'emerge pr\'ecis\'ement de cette dualit\'e, puisqu'une preuve peut
\^etre cryptographiquement impeccable ($\mathsf{Verify}(\sigma) = 1$) mais
institutionnellement opaque, l'interpr\`ete \'etant incapable d'en extraire
$v$ de mani\`ere fiable.
\end{remark}

\begin{figure}[htbp]
\centering
\begin{tikzpicture}[
    scale=0.85,
    transform shape,
    box/.style={rectangle, draw=enspblue, fill=blue!8, thick,
                minimum width=2.8cm, minimum height=1cm,
                align=center, font=\small},
    cryptobox/.style={rectangle, draw=enspgreen, fill=enspgreen!8, thick,
                minimum width=2.8cm, minimum height=1cm,
                align=center, font=\small},
    instbox/.style={rectangle, draw=navyblue, fill=navyblue!8, thick,
                minimum width=2.8cm, minimum height=1cm,
                align=center, font=\small},
    dataarrow/.style={-{Stealth[length=3mm]}, thick, enspblue},
    interpretarrow/.style={-{Stealth[length=3mm]}, thick, navyblue, dashed}
]

% ---------------------------------------------------------------
% N\OE UDS
% ---------------------------------------------------------------

% Ligne 1 : entr\'ees
\node[box] (message) at (0, 6)   {\textbf{Message} $M$};
\node[box] (context) at (6, 6)   {\textbf{Contexte} $C$};

% Ligne 2 : Prouveur
\node[cryptobox, minimum width=4cm] (prover) at (3, 4)
    {\textbf{Prouveur}\\[2pt]
     $\mathsf{Prouver}(sk, m, c) \to \sigma$};

% Ligne 3 : Preuve
\node[cryptobox] (proof) at (3, 2)
    {\textbf{Preuve} $\Sigma$};

% Ligne 4 gauche : V\'erifier
\node[cryptobox] (verify) at (0, 0)
    {\textbf{V\'erifier}\\[2pt]
     $\mathsf{V\'erifier}(pk, \sigma, c)$};

% Ligne 5 gauche : r\'esultat binaire
\node[box] (vresult) at (0, -1.8)
    {$\{0,1\}$};

% Ligne 4 droite : Interpr\'eter
\node[instbox] (interpret) at (7, 0)
    {\textbf{Interpr\'eter}\\[2pt]
     $\mathsf{Interpr\'eter}(\sigma, c)$};

% Ligne 5 droite : Verdict
\node[instbox] (verdict) at (7, -1.8)
    {\textbf{Verdict} $V$};

% ---------------------------------------------------------------
% FL\`ECHES COUCHE CRYPTOGRAPHIQUE
% ---------------------------------------------------------------
\draw[dataarrow] (message.south) -- (prover.north west);
\draw[dataarrow] (context.south) -- (prover.north east);
\draw[dataarrow] (prover.south)  -- (proof.north);
\draw[dataarrow] (proof.south)   -- (verify.north);
\draw[dataarrow] (context.west)  -| (verify.north);
\draw[-{Stealth[length=3mm]}, thick]
                 (verify.south)  -- (vresult.north);

% ---------------------------------------------------------------
% FL\`ECHES COUCHE INSTITUTIONNELLE
% ---------------------------------------------------------------
\draw[interpretarrow] (proof.east)    -| (interpret.north);
\draw[interpretarrow] (context.east)  -| (interpret.north);
\draw[interpretarrow] (interpret.south) -- (verdict.north);

% ---------------------------------------------------------------
% \'ETIQUETTES FLOTTANTES (sans fit, pour \'eviter les chevauchements)
% ---------------------------------------------------------------
\node[font=\footnotesize, text=enspblue]  at (-1.2, 5.4)
    {Entr\'ee};
\node[font=\footnotesize, text=enspgreen] at (-1.4, 3.0)
    {Cryptographique};
\node[font=\footnotesize, text=enspgreen] at (1.5,  2.5)
    {Sortie};
\node[font=\footnotesize, text=navyblue]  at (8.6,  1.0)
    {Institutionnel};

% ---------------------------------------------------------------
% CADRES DE COUCHES --- dessin\'es MANUELLEMENT pour contr\^ole pr\'ecis
% ---------------------------------------------------------------

% Cadre couche cryptographique (englobe prover, proof, verify, vresult)
\draw[enspgreen, dashed, thick, rounded corners=4pt]
    (-1.6, -2.5) rectangle (5.0, 4.8);
\node[font=\small\bfseries, text=enspgreen] at (1.7, 4.5)
    {Couche cryptographique};

% Cadre couche institutionnelle (englobe interpret, verdict)
\draw[navyblue, dashed, thick, rounded corners=4pt]
    (5.4, -2.5) rectangle (9.0, 0.8);
\node[font=\small\bfseries, text=navyblue] at (7.2, 0.55)
    {Couche institutionnelle};

\end{tikzpicture}
\caption[Mod\`ele de Protocole CLO : architecture de v\'erification \`a deux couches]{%
  Mod\`ele de Protocole CLO : architecture de v\'erification \`a deux couches.%
}\label{fig:clo_model}

\smallskip
\begin{minipage}{\linewidth}
  \small\textit{Note :} La couche cryptographique (verte) assure la correction
  math\'ematique via les algorithmes $\mathsf{Prouver}$ et $\mathsf{V\'erifier}$.
  La couche institutionnelle (bleue marine) fournit l'interpr\'etation s\'emantique
  via la fonction $\mathsf{Interpr\'eter}$, qui projette les preuves cryptographiques
  sur des verdicts juridiquement significatifs en contexte institutionnel.
\end{minipage}

\end{figure}
\begin{boiteArbitre}
L'arbitre comprend que l'algorithme $\mathsf{Interpret}$ est le sien, car les trois
autres prot\`egent le message alors que celui-ci le sert, lui.
\end{boiteArbitre}

\subsubsection{Exemples d'instanciation}
\label{subsubsec:exemples_instanciation}

Les exemples qui suivent ont valeur illustrative~; les m\'etriques CLO
num\'eriques qui leur seront associ\'ees dans la suite sont des estimations
param\'etriques analytiques d\'eriv\'ees depuis les sp\'ecifications publi\'ees,
non des mesures empiriques contr\^ol\'ees.

\paragraph{Vote \'electronique (Helios)~\cite{Adida08,Adida09UCL}.}
Le message $m$ est le choix de vote $\{0,1\}$~; la preuve $\sigma$
consiste en un bulletin chiffr\'e ElGamal accompagn\'e d'une preuve NIZK~(Non-Interactive Zero-Knowledge)
de bien-formation~\cite{Chaum81,CramerDamgardSchoenmakers}. Le contexte $c$ recouvre le registre d'\'electeurs,
la structure du bulletin et les r\`egles \'electorales, tandis que le
verdict $v$ d\'esigne le d\'ecompte final ou la validation de l'\'election.
$\mathsf{Verify}$ contr\^ole la preuve NIZK et la signature du bulletin~;
$\mathsf{Interpret}$ consiste pour les garants \'electoraux \`a d\'echiffrer
l'agr\'egat homomorphe et en extraire le d\'ecompte~\cite{Goodman15}. Les m\'etriques CLO
estim\'ees (\'echelle op\'erationnelle CLO de~\cite{MINKA2026CLO}, fiabilit\'e
$\Relver$) sont $\Conf = 0{,}92$, $\Relver = 0{,}95$, $\Opp = 0{,}50$, avec
un produit $\Conf\cdot\Relver\cdot\Opp = 0{,}437$.

\paragraph{Conformit\'e financi\`ere (zk-SNARKs)~\cite{Groth16,BCTV14,Gennaro13,Parno13}.}
Ici $m$ encode transactions priv\'ees ou soldes~; $\sigma$ est une preuve
Groth16 de 128 octets certifiant la conformit\'e r\'eglementaire~\cite{BenSasson2019,PLONK20}. Le contexte
$c$ couvre les r\'egulations AML/B\^ale~III~\cite{BaselIII}, l'historique de
transactions et les politiques internes. $\mathsf{Verify}$ contr\^ole les
\'equations de \emph{pairing} de Groth16~\cite{Campanelli2019,Wahby18}~; $\mathsf{Interpret}$ demande au
r\'egulateur de corr\'eler le hachage de contexte avec ses bases de donn\'ees
pour rendre un verdict de conformit\'e~\cite{FATF15}. Les m\'etriques CLO estim\'ees sont
$\Conf = 0{,}89$, $\Relver = 0{,}96$, $\Opp = 0{,}67$, avec un produit
$\Conf\cdot\Relver\cdot\Opp = 0{,}572$.

\paragraph{Signatures contractuelles (ECDSA)~\cite{Johnson01,NIST.FIPS.186-5}.}
Le message $m$ est le document contractuel~; $\sigma = (r, s)$ est la
signature ECDSA. Le contexte $c$ comprend les parties contractantes, la
juridiction et les termes du contrat~; $v$ est la validit\'e du contrat et
l'identification de la partie responsable~\cite{Mason07,Kerr04}. $\mathsf{Verify}$ contr\^ole
l'\'equation $r = x_1 \bmod n$~; $\mathsf{Interpret}$ engage le juge dans
la v\'erification des identit\'es, l'examen des termes et le rendu du jugement.
Les m\'etriques CLO estim\'ees sont $\Conf = 0{,}31$, $\Relver = 0{,}98$, $\Opp = 0{,}90$,
avec un produit $\Conf\cdot\Relver\cdot\Opp = 0{,}273$.

Ces trois instanciations couvrent vote \'electronique, r\'egulation financi\`ere
et litige contractuel~\cite{Nakamoto08,Buterin17}, registre suffisant pour tester la g\'en\'eralit\'e du
mod\`ele. Le tableau~\ref{tab:clo_resume_ch1} r\'ecapitule ces m\'etriques. Ces
syst\`emes ne sont pas des instances du Trilemme CRO
(Remarque~\ref{rem:portee_cas_etude})~; ils servent \`a calibrer les ordres
de grandeur des param\`etres institutionnels.

% -----------------------------------------------------------------------
%\begin{table}[htbp]
%\centering
%\caption{M\'etriques CLO des syst\`emes d\'eploy\'es}
%\label{tab:clo_resume_ch2_clo}
%\begin{tabular}{@{}lcccc@{}}
%\toprule
%\textbf{Syst\`eme} & $\boldsymbol{\Conf}$ & $\boldsymbol{\Rel}$ & $\boldsymbol{\Opp}$ & $\boldsymbol{\Conf\cdot\Rel\cdot\Opp}$ \\
%\midrule
%Helios (vote \'electronique)~\cite{Adida08,Adida09UCL} & 0{,}92 & 0{,}95 & 0{,}50 & \textbf{0{,}437} \\
%zk-SNARKs (conformit\'e)~\cite{Groth16,BCTV14} & 0{,}89 & 0{,}96 & 0{,}67 & \textbf{0{,}572} \\
%ECDSA (contrats)~\cite{Johnson01,NIST.FIPS.186-5} & 0{,}31 & 0{,}98 & 0{,}90 & \textbf{0{,}273} \\
%\bottomrule
%\end{tabular}
%\begin{tablenotes}
%\item \footnotesize Source~: analyses param\'etriques d\'eriv\'ees des sp\'ecifications publi\'ees~\cite{Tsybakov2009}.
%\item \footnotesize Ces valeurs servent de validation th\'eorique, non de mesures empiriques.
%\end{tablenotes}
%\end{table}
\end{sloppypar}

% -----------------------------------------------------------------------
\subsection{Les trois dimensions : d\'efinitions formelles}
\label{subsec:trois_dimensions}
% -----------------------------------------------------------------------
\begin{sloppypar}
Nous d\'efinissons les trois m\'etriques fondamentales via la th\'eorie
de l'information (section~\ref{sec:etat_art},
sous-section~\ref{subsec:fondements_info})~\cite{Shannon48,CoverThomas06,HanVerdu}.

\subsubsection{Confidentialit\'e : indistinguabilit\'e au pire cas}
\label{subsubsec:confidentialite_clo}

La confidentialit\'e formalise l'impossibilit\'e de distinguer les preuves
g\'en\'er\'ees pour des messages distincts, g\'en\'eralisant l'indistinguabilit\'e
s\'emantique~\cite{GoldwasserMicali84}.

\begin{definition}[Confidentialit\'e]\label{def:confidentialite_clo}
Pour un protocole CLO $\mathcal{P}$ :
\begin{equation}
\Conf(\mathcal{P}) = 1 - \max_{\substack{m_0, m_1 \in \mathcal{M}_\lambda \\ c \in \mathcal{C}_\lambda}}
\TV\!\left(P_{\Sigma|M=m_0,C=c},\; P_{\Sigma|M=m_1,C=c}\right)
\label{eq:conf_clo}
\end{equation}
o\`u $\TV(P, Q) = \frac{1}{2} \sum_{\sigma \in \Sigma_\lambda} |P(\sigma) - Q(\sigma)|$
est la distance en variation totale.
\end{definition}

$\Conf = 1$ (confidentialit\'e parfaite au sens de Shannon~\cite{Shannon49,Shannon48})
signifie que les distributions de preuves sont identiques pour tous
messages et tout contexte~; m\^eme un adversaire computationnellement
illimit\'e ne peut distinguer~\cite{Wyner75,WynerZiv76}. \`a l'oppos\'e, $\Conf = 0$ trahit l'existence
de messages $m_0, m_1$ et d'un contexte $c$ pour lesquels les supports
sont disjoints, un observateur passif acquiert alors une certitude.
Le r\'egime pratique est interm\'ediaire.

Du point de vue des jeux de s\'ecurit\'e, notre d\'efinition majore l'avantage
de tout distingueur dans IND-CPA~\cite{GoldwasserMicali84,Bellare97}~:
\begin{equation}
\mathsf{Adv}^{\text{IND-CPA}}_{\mathcal{D}}(\lambda)
  \leq \max_{m_0,m_1,c} \TV\!\left(P_{\Sigma|m_0,c},\; P_{\Sigma|m_1,c}\right)
\label{eq:adv_ind_cpa}
\end{equation}
$\Conf$ fournit donc une borne informationnelle sur la s\'ecurit\'e cryptographique.

\subsubsection{Fiabilit\'e : reconstruction du message}
\label{subsubsec:fiabilite_clo}

\begin{definition}[Fiabilit\'e de reconstruction]\label{def:fiabilite_clo}
Soit $\hat{M} = \mathsf{Dec}(\Sigma, C)$ un estimateur de $M$ calculable \`a
partir de la preuve et du contexte. La fiabilit\'e est
\begin{equation}
\Rel(\mathcal{P}) \;:=\; \Relrec(\mathcal{P})
  = 1 - \min_{\mathsf{Dec}} \Pr\bigl[\mathsf{Dec}(\Sigma, C) \neq M\bigr],
\label{eq:rel_clo}
\end{equation}
la probabilit\'e \'etant prise sur $M$, $C$ et l'al\'ea du protocole
(\'equation~\eqref{eq:rel_def}).
\end{definition}

La fiabilit\'e mesure la capacit\'e de la preuve \`a restituer le message
original~\cite{Fano61,CoverThomas06}. C'est cette seule notion que
contr\^ole l'in\'egalit\'e de Fano. Un \emph{one-time pad} pur atteint
$\Conf = 1$ mais $\Rel = 2^{-\lambda}$ pour $M$ uniforme (pas de
r\'ecup\'eration sans cl\'e)~\cite{Shannon49}~; un engagement de
Pedersen~\cite{Pedersen91,ChaumPedersen92} est parfaitement masquant
($\Conf = 1$, donc $\Relrec = 2^{-\lambda}$) et calculatoirement liant sous
DLog~\cite{Damgard87}~; ECDSA~\cite{Johnson01,NIST.FIPS.186-5}, qui publie le message
en clair, a $\Relrec = 1$ et une confidentialit\'e nulle au sens de la
D\'efinition~\ref{def:confidentialite_clo}.

\begin{remark}[Fiabilit\'e de v\'erification]\label{rem:rel_ver}
La fiabilit\'e \emph{op\'erationnelle} de~\cite{MINKA2026CLO} est la
compl\'etude, assortie de la solidit\'e, de la v\'erification~:
\begin{equation}
\Relver(\mathcal{P}) = \min_{m,c}\Pr\bigl[\mathsf{Verify}(pk,\mathsf{Prove}(sk,m,c),c)=1\bigr].
\label{eq:rel_ver}
\end{equation}
Les deux notions ne co\"incident pas. Pour $M$ uniforme sur $\{0,1\}^\lambda$,
$\Conf = 1$ entra\^ine $\Sigma \perp M \mid C$, donc $\Relrec = 2^{-\lambda}$,
alors que $\Relver$ peut valoir~$1$ (engagement parfaitement masquant). Le
Trilemme (Th\'eor\`eme~\ref{thm:trilemme_cro_conditionnel}) porte sur
$\Relrec$, seule notion \`a laquelle s'applique l'in\'egalit\'e de Fano. Les
m\'etriques CLO des syst\`emes d\'eploy\'es et des constructions SOTP sont
exprim\'ees avec $\Relver$~; elles ne sont donc pas des instances du Trilemme
(Remarque~\ref{rem:portee_cas_etude}).
\end{remark}

\subsubsection{Opposabilit\'e : r\'esolution d'incertitude institutionnelle}
\label{subsubsec:opposabilite_clo}

\begin{definition}[Opposabilit\'e]\label{def:opposabilite_clo}
\begin{equation}
\Opp(\mathcal{P}) = \frac{I(V;\Sigma|C)}{H(V|C)}
\label{eq:opp_clo}
\end{equation}
\end{definition}

L'opposabilit\'e capture la fraction de l'incertitude du verdict
institutionnel $V$ que la preuve $\Sigma$ r\'esout, conditionnellement
au contexte $C$. Absente de la litt\'erature cryptographique
classique, cette m\'etrique quantifie
l'interpr\'etabilit\'e institutionnelle. Ainsi, $\Opp = 1$ signifie que la preuve
l\`eve toute ambigu\"it\'e~; $\Opp = 0$ que la preuve est institutionnellement
muette~\cite{Zohar15,Grimmelmann15}. L'information mutuelle conditionnelle $I(V;\Sigma|C)$ est born\'ee
par la taille de l'espace de verdicts, point que le th\'eor\`eme~\ref{thm:goulot_contexte}
exploitera. La d\'efinition n'est informative que si le verdict n'est pas une
fonction d\'eterministe de $(\Sigma,C)$, car dans ce cas $H(V|\Sigma,C)=0$ et
$\Opp=1$ pour tout protocole (Remarque~\ref{rem:verdict_deterministe}). Le
cadre retient donc des interpr\'etations admissibles au sens de la
D\'efinition~\ref{def:interp_admissible}, o\`u le verdict d\'epend aussi des
ressources propres de l'institution.

\end{sloppypar}

% -----------------------------------------------------------------------
\subsection{Propri\'et\'es fondamentales}
\label{subsec:proprietes_fondamentales}
% -----------------------------------------------------------------------
\begin{sloppypar}
Nous \'etablissons quatre propri\'et\'es structurant les m\'etriques CLO.
Les deux premi\`eres sont des bornes \'el\'ementaires~; les deux suivantes
r\'ev\`elent des tensions intrins\`eques.

\begin{theorem}[Bornes \'el\'ementaires des m\'etriques CLO]\label{thm:bornes_mesures}
Pour tout protocole CLO $\mathcal{P}$~:
\begin{equation}
  0 \leq \Conf(\mathcal{P}),\; \Rel(\mathcal{P}),\; \Opp(\mathcal{P}) \leq 1
  \label{eq:bornes_elementaires}
\end{equation}
\end{theorem}

\begin{Preuve}[Th\'eor\`eme~\ref{thm:bornes_mesures}~: Bornes \'el\'ementaires]\label{prv:bornes_mesures}

\begin{cmtEtape}{\'Etape unique. Montrer que chaque m\'etrique est dans $[0,1]$.}
Les trois m\'etriques sont d\'efinies comme des rapports ou des
compl\'ements de distances. On v\'erifie s\'epar\'ement que chacune
reste dans l'intervalle $[0,1]$.
\end{cmtEtape}

\textbf{Confidentialit\'e.}
Par d\'efinition, $\Conf(\mathcal{P}) = 1 - \max_{m_0,m_1}
\TV(P_{\Sigma|M=m_0}, P_{\Sigma|M=m_1})$.
\cmt{La confidentialit\'e est l'oppos\'e de la meilleure distinguabilit\'e.}

Puisque $\TV(P,Q) \in [0,1]$ pour toutes distributions $P$ et $Q$,
\cmt{La distance en variation totale est comprise entre 0 et 1.}

on a $0 \leq 1 - \max \TV(\cdot) \leq 1$,
\cmt{Si le max vaut 0, $\Conf = 1$. Si le max vaut 1, $\Conf = 0$.}

d'o\`u $\Conf \in [0,1]$.
\cmt{Conclusion pour la confidentialit\'e.}

\textbf{Fiabilit\'e.}
Par d\'efinition, $\Rel(\mathcal{P}) = 1 - \min_{\mathsf{Dec}} \Pr[\mathsf{Dec}(\Sigma,C) \neq M]$.
\cmt{La fiabilit\'e est la meilleure probabilit\'e de reconstruction du message.}

Une probabilit\'e d'erreur est dans $[0,1]$, d'o\`u $\Rel \in [0,1]$~; il en va
de m\^eme pour $\Relver$ (Remarque~\ref{rem:rel_ver}).
\cmt{Une probabilit\'e est toujours dans $[0,1]$.}

\textbf{Opposabilit\'e.}
Par d\'efinition, $\Opp(\mathcal{P}) = I(V;\Sigma|C)/H(V|C)$.
\cmt{Rapport de l'information extraite sur l'information disponible.}

Par les propri\'et\'es fondamentales de l'information mutuelle
conditionnelle~\cite{CoverThomas06},
\begin{equation}
0 \leq I(V;\Sigma|C) \leq H(V|C)
\end{equation}
\cmtFormule{L'information mutuelle est toujours born\'ee par les entropies marginales.}

ce qui donne $\Opp \in [0,1]$ en divisant par $H(V|C) > 0$.
\cmt{Le rapport de deux quantit\'es positives est dans $[0,1]$ si le num\'erateur est major\'e par le d\'enominateur.}
\end{Preuve}
\begin{lemma}[Tension confidentialit\'e--information mutuelle]\label{lem:compromis_conf_mi}
Pour les protocoles $\beta$-r\'eguliers (Hypoth\`ese~\ref{ass:regularite})~:
\begin{equation}
\Conf \geq 1 - \sqrt{c_0\,\beta \cdot I(M;\Sigma|C)}
\label{eq:compromis_conf_mi}
\end{equation}
o\`u $c_0 = 1/(2\ln 2) \approx 0{,}721$. Cette minoration est la forme
utilis\'ee dans toute la th\`ese~; elle est reprise, sous sa forme
d\'efinitive et avec ses hypoth\`eses compl\`etes, au
Lemme~\ref{lem:confidentiality_guarantee}.
\end{lemma}

\begin{Preuve}[Lemme~\ref{lem:compromis_conf_mi}~: Trade-off Conf-MI]\label{prv:compromis_conf_mi}

\begin{cmtEtape}{Objectif.}
\'Etablir que $\Conf \geq 1 - \sqrt{c_0 \beta I}$. Autrement dit, tant
que la preuve divulgue peu d'information sur le message, la
confidentialit\'e reste garantie~; c'est lorsque $I$ cro\^it que cette
garantie s'affaiblit. C'est la premi\`ere moiti\'e du m\'ecanisme du
trilemme.
\end{cmtEtape}

Par l'Hypoth\`ese~\ref{ass:regularite} ($\beta$-r\'egularit\'e) et l'in\'egalit\'e de
Pinsker (Lemme~\ref{lem:pinsker}), pour tout couple $(m_0, m_1)$ et
contexte $c$~:
\begin{align}
\TV^2(P_{\Sigma|m_0,c}, P_{\Sigma|m_1,c})
&\leq c_0 \KL(P_{\Sigma|m_0,c} \| P_{\Sigma|m_1,c})
\\
&\leq c_0 \beta \cdot I(M;\Sigma|C)
\end{align}
\cmtFormule{Application directe de Pinsker.}
\cmtFormule{Application de l'hypoth\`ese de $\beta$-r\'egularit\'e.}
\cmt{La distance TV est born\'ee par la divergence KL.}
\cmt{La divergence KL est born\'ee par $\beta \cdot I(M;\Sigma|C)$.}

En prenant le maximum sur tous les couples et contextes, on obtient
\begin{equation}
\max_{\substack{m_0, m_1 \in \mathcal{M}_\lambda \\ c \in \mathcal{C}_\lambda}}
\TV(P_{\Sigma|m_0,c}, P_{\Sigma|m_1,c})
\leq \sqrt{c_0\beta \cdot I(M;\Sigma|C)}
\end{equation}
\cmtFormule{Le pire cas de distinguabilit\'e est born\'e par $\sqrt{c_0 \beta I}$.}

Puisque $\Conf(\mathcal{P}) = 1 - \max \TV(\cdots)$ par d\'efinition
de $\Conf$, il vient
\begin{equation}
\Conf(\mathcal{P}) \geq 1 - \sqrt{c_0\,\beta \cdot I(M;\Sigma|C)}
\end{equation}
\cmtFormule{La garantie de confidentialit\'e d\'ecro\^it en racine de l'information mutuelle.}
\cmt{Aucune normalisation par $\lambda$ n'intervient, car l'argument de Pinsker porte sur des lois conditionnelles, non sur la taille de l'espace de messages.}
\end{Preuve}

\begin{remark}[Interpr\'etation du lemme]\label{rem:interpretation_lemme}
Ce lemme fournit une \emph{garantie} de confidentialit\'e. Plus
$I(M;\Sigma|C)$ est petit, plus le minorant de $\Conf$ est proche de~$1$,
selon une d\'ependance non lin\'eaire (en racine carr\'ee). La tension
$\Conf$--$\Rel$ elle-m\^eme vient de l'in\'egalit\'e de Fano, puisqu'une fiabilit\'e
de reconstruction \'elev\'ee exige que la preuve contienne de l'information
sur le message. Un minorant de $\Conf$ ne suffit pas, \`a lui seul, \`a
formaliser cette tension.
\end{remark}

\begin{lemma}[Lemme de Johnson-Lindenstrauss~\cite{JohnsonLindenstrauss84}]\label{lem:johnson_lindenstrauss}
Pour tout $0 < \varepsilon < 1$, tout entier $n > 0$ et tout ensemble 
 $S \subset \mathbb{R}^n$ avec $|S| = N$, il existe une application lin\'eaire 
 $\Phi: \mathbb{R}^n \to \mathbb{R}^d$ avec $d = O(\varepsilon^{-2} \log N)$ 
telle que pour tous $x, y \in S$ :
\begin{equation}
(1-\varepsilon)\|x-y\|^2 \leq \|\Phi x - \Phi y\|^2 \leq (1+\varepsilon)\|x-y\|^2
\label{eq:jl_preservation}
\end{equation}
avec haute probabilit\'e lorsque $\Phi$ est une matrice al\'eatoire avec des 
entr\'ees $\Phi_{ij} \sim \mathcal{N}(0, 1/d)$ ou sous-gaussiennes.
\end{lemma}

Ce lemme garantit la pr\'eservation des distances euclidiennes sous compression
al\'eatoire, propri\'et\'e sur laquelle repose la construction SOTP-Base
(section~\ref{sec:sotp_construction}, sous-section~\ref{subsec:sotp_base}),
puisque les distances entre repr\'esentations de contextes sont pr\'eserv\'ees \`a un
facteur $1 \pm \varepsilon$ pr\`es apr\`es compression. Que l'interpr\`ete
institutionnel puisse en extraire autant d'information d\'ecisionnelle
qu'avant compression est une hypoth\`ese de conception, non une cons\'equence
du lemme.

\begin{theorem}[Goulot d'\'etranglement contextuel]\label{thm:goulot_contexte}
Pour tout protocole CLO~:
\begin{equation}
I(V;\Sigma|C) \leq \min\bigl\{H(\Sigma|C),\; H(V|C)\bigr\}
  \leq \log_2|\mathcal{V}_\lambda| \leq \alpha\, H(\mathcal{C})
\label{eq:goulot_contextuel}
\end{equation}
\end{theorem}

\begin{Preuve}[Th\'eor\`eme~\ref{thm:goulot_contexte}~: Goulot d'\'etranglement]\label{prv:goulot_contexte}

\begin{cmtEtape}{Objectif.}
Montrer que $I(V;\Sigma|C) \leq \min\{H(\Sigma|C),\, H(V|C)\} \leq \alpha H(C)$.
Autrement dit, la preuve $\Sigma$ ne peut transmettre plus d'information
sur le verdict que ce qu'elle contient, ni plus que l'incertitude
initiale sur le verdict, laquelle est plafonn\'ee par la granularit\'e des
verdicts rapport\'ee \`a la richesse du contexte.
\end{cmtEtape}

\textbf{Borne 1 (contrainte par la capacit\'e de preuve).}
Par d\'efinition de l'information mutuelle conditionnelle, on a
\begin{equation}
I(V;\Sigma|C) = H(\Sigma|C) - H(\Sigma|V,\mathcal{C})
\end{equation}
\cmtFormule{L'information sur $V$ contenue dans $\Sigma$ est l'entropie de $\Sigma$ moins ce qui reste une fois $V$ connu.}

Puisque l'entropie conditionnelle est non-n\'egative~\cite{CoverThomas06},
$H(\Sigma|V,\mathcal{C}) \geq 0$, on obtient
\begin{equation}
I(V;\Sigma|C) \leq H(\Sigma|C)
\end{equation}
\cmtFormule{Enlever un terme positif donne une in\'egalit\'e.}
\cmt{On ne peut pas extraire plus d'information que la preuve n'en contient.}

\textbf{Borne 2 (contrainte par l'incertitude sur les verdicts).}
Par la propri\'et\'e sym\'etrique de l'information mutuelle, on a
\begin{equation}
I(V;\Sigma|C) = H(V|C) - H(V|\Sigma,\mathcal{C})
\end{equation}
\cmtFormule{Formulation sym\'etrique, o\`u l'information \'egale l'incertitude initiale moins l'incertitude r\'esiduelle.}

Puisque $H(V|\Sigma,\mathcal{C}) \geq 0$, il vient
\begin{equation}
I(V;\Sigma|C) \leq H(V|C)
\end{equation}
\cmtFormule{Le verdict ne peut \^etre r\'esolu au-del\`a de son incertitude initiale.}
\cmt{On ne peut pas r\'esoudre plus d'incertitude qu'il n'y en a.}

\textbf{Borne sur $H(V|C)$.}
L'entropie d'une variable est major\'ee par le logarithme de la taille de
son support~\cite{CoverThomas06}. Or la granularit\'e des verdicts
(D\'efinition~\ref{def:alpha_granularite},
\'equation~\eqref{eq:alpha_definition}) borne ce support, d'o\`u
\begin{equation}
H(V|C) \leq \log_2 |\mathcal{V}_\lambda| \leq \alpha\, H(\mathcal{C})
\end{equation}
\cmtFormule{L'incertitude sur le verdict ne peut d\'epasser le nombre de verdicts distincts.}
\cmt{C'est la granularit\'e $\alpha$, param\`etre institutionnel, qui relie le plafond au contexte.}

\textbf{Conclusion.}
La combinaison des trois bornes donne
\begin{equation}
I(V;\Sigma|C) \leq \min\{H(\Sigma|C),\; H(V|C)\} \leq \alpha\, H(C)
\end{equation}
\cmtFormule{Le plafond de l'extraction est fix\'e par la granularit\'e institutionnelle des verdicts.}
\end{Preuve}

\begin{remark}[Verdict d\'eterministe]\label{rem:verdict_deterministe}
Si $V$ est une fonction d\'eterministe de $(\Sigma,C)$, alors
$H(V|\Sigma,C)=0$ et $I(V;\Sigma|C)=H(V|C)$, de sorte que l'opposabilit\'e vaut $1$
d\`es que $H(V|C)>0$, quel que soit le protocole. La m\'etrique $\Opp$ ne
discrimine les protocoles que pour des interpr\'etations admissibles au sens
de la D\'efinition~\ref{def:interp_admissible}, o\`u le verdict d\'epend
aussi de ressources institutionnelles $K$ non r\'eductibles \`a $(\Sigma,C)$.
Ce fait est certifi\'e en Lean (\texttt{opp\_eq\_one\_of\_deterministic},
section~\ref{sec:lean4}).
\end{remark}

Ce th\'eor\`eme \'etablit qu'aucune sophistication cryptographique ne saurait
d\'epasser $\alpha H(\mathcal{C})$ bits d'information extractible par un interpr\`ete
institutionnel~\cite{Smith09,Alvim12}. La granularit\'e des verdicts rapport\'ee \`a
l'entropie contextuelle constitue un plafond sur
l'opposabilit\'e, plafond qui pr\'efigure le Trilemme~CRO classique de la section~\ref{sec:cro_classique} et son extension quantique de la section~\ref{sec:cro_pq}. Les protocoles dans des contextes
simples ($H(\mathcal{C}) \approx 20$ bits, comme ECDSA pour contrat
standardis\'e) peuvent atteindre haute opposabilit\'e sans effort~; les
contextes complexes ($H(\mathcal{C}) \approx 70$ bits, comme zk-SNARKs
pour conformit\'e AML) imposent un plafond sauf si la capacit\'e
institutionnelle $\Delta$ est exceptionnellement \'elev\'ee.

\begin{boiteArbitre}
L'arbitre comprend que le plafond tient moins \`a la cryptographie qu'\`a
la finesse de ses propres verdicts, rapport\'ee au contexte. Plus le contexte est simple, plus il peut
d\'ecider, sans que la preuve y puisse quoi que ce soit.
\end{boiteArbitre}

\begin{theorem}[R\`egles de composition]\label{thm:composition_sequentielle}
Soit $\mathcal{P}_1 \circ \mathcal{P}_2$ le protocole qui publie
$\Sigma_2 = \mathcal{P}_2(\Sigma_1)$, o\`u $\Sigma_1$ est produite par
$\mathcal{P}_1$ (post-traitement, \'eventuellement randomis\'e,
ind\'ependant de $M$ sachant $(\Sigma_1, C)$). Alors
\begin{align}
\Conf(\mathcal{P}_1 \circ \mathcal{P}_2) &\geq \Conf(\mathcal{P}_1),
  \label{eq:comp_seq_conf}\\
\Relrec(\mathcal{P}_1 \circ \mathcal{P}_2) &\leq \Relrec(\mathcal{P}_1).
  \label{eq:comp_seq_rel}
\end{align}
Soit $\mathcal{P}_1 \| \mathcal{P}_2$ le protocole qui publie
$(\Sigma_1,\Sigma_2)$ produits ind\'ependamment sachant $(M,C)$. Alors
\begin{align}
\Conf(\mathcal{P}_1 \| \mathcal{P}_2) &\leq \min\bigl\{\Conf(\mathcal{P}_1),\, \Conf(\mathcal{P}_2)\bigr\},
  \label{eq:comp_par_conf}\\
\Relrec(\mathcal{P}_1 \| \mathcal{P}_2) &\geq \max\bigl\{\Relrec(\mathcal{P}_1),\, \Relrec(\mathcal{P}_2)\bigr\}.
  \label{eq:comp_par_rel}
\end{align}
\end{theorem}

\begin{Preuve}[Th\'eor\`eme~\ref{thm:composition_sequentielle}~: R\`egles de composition]\label{prv:composition}

\begin{cmtEtape}{Objectif.}
\'Etablir le comportement de la confidentialit\'e et de la fiabilit\'e sous
post-traitement (composition s\'equentielle) et sous publication conjointe
(composition parall\`ele).
\end{cmtEtape}

\textbf{Composition s\'equentielle.}
$M \to (\Sigma_1,C) \to (\Sigma_2,C)$ est une cha\^ine de Markov.
L'in\'egalit\'e de traitement des donn\'ees pour la distance en variation
totale~\cite{CoverThomas06} donne, pour tous $m_0, m_1, c$~:
\begin{equation}
\TV(P_{\Sigma_2|m_0,c}, P_{\Sigma_2|m_1,c}) \leq \TV(P_{\Sigma_1|m_0,c}, P_{\Sigma_1|m_1,c}),
\end{equation}
\cmtFormule{Un traitement ne cr\'ee pas de distinguabilit\'e.}
d'o\`u~\eqref{eq:comp_seq_conf}. Tout d\'ecodeur de $M$ \`a partir de
$(\Sigma_2,C)$ en induit un \`a partir de $(\Sigma_1,C)$, obtenu en appliquant d'abord
$\mathcal{P}_2$ \`a $\Sigma_1$, puis ce d\'ecodeur. D'o\`u~\eqref{eq:comp_seq_rel}.
\cmt{Post-traiter une preuve ne peut que la rendre plus confidentielle et moins reconstructible.}

\textbf{Composition parall\`ele.}
La marginale de $(\Sigma_1,\Sigma_2)$ sur $\Sigma_i$ est un post-traitement
de la paire. Par la m\^eme in\'egalit\'e, la distance en variation totale
entre lois conjointes majore celle entre marginales, pour chaque $i$~;
d'o\`u~\eqref{eq:comp_par_conf}. Un d\'ecodeur qui ignore l'une des
composantes est admissible pour la paire, d'o\`u~\eqref{eq:comp_par_rel}.
\cmt{Publier deux preuves r\'ev\`ele au moins autant que chacune d'elles.}
\end{Preuve}

\begin{remark}[Opposabilit\'e compos\'ee]\label{rem:composition_opp}
Pour l'opposabilit\'e, la cha\^ine $V$--$(\Sigma_1,C)$--$(\Sigma_2,C)$ donne
$I(V;\Sigma_2|C) \leq I(V;\Sigma_1|C)$, donc
$\Opp(\mathcal{P}_1\circ\mathcal{P}_2) \leq \Opp(\mathcal{P}_1)$ \`a verdict
fix\'e. Aucune r\`egle g\'en\'erale ne vaut pour la composition parall\`ele,
car l'information mutuelle n'est ni sous- ni sur-additive en g\'en\'eral. Les
r\`egles de~\cite{MINKA2026CLO} ($\Opp \geq \min$ en s\'equentiel,
$\Opp \approx$ produit en parall\`ele) sont des heuristiques qui ne sont pas
reprises ici. La formulation initiale de ce th\'eor\`eme dans~\cite{MINKA2026CLO}
(minoration par le produit via l'in\'egalit\'e de Fr\'echet, $\Conf \geq \max$
en parall\`ele) est remplac\'ee ici par une formulation plus pr\'ecise,
l'in\'egalit\'e de Fr\'echet donnant $ab \geq a+b-1$.
\end{remark}

Ces r\`egles permettent d'analyser des protocoles composites. Ainsi, un syst\`eme
bancaire publiant conjointement une signature ECDSA (authentification) et une preuve zk-SNARK
(conformit\'e) h\'erite des bornes ci-dessus, puisque publier les deux preuves ne peut
qu'abaisser la confidentialit\'e par rapport \`a chacune d'elles.

\subsubsection{Jeux cryptographiques et m\'etriques CLO}
\label{subsubsec:jeux_crypto}

Les d\'efinitions informationnelles de $\Conf$, $\Rel$ et $\Opp$ sont
compl\'et\'ees par des jeux cryptographiques qui captent leur contenu
op\'erationnel en termes d'avantages d'adversaires PPT. Cette section
\'etablit le pont formel entre les deux formalismes.

\paragraph{Jeu IND-CONF : confidentialit\'e game-based.}
\label{subsubsubsec:jeu_ind_conf}

\begin{definition}[Jeu IND-CONF]\label{def:jeu_ind_conf}
Pour un protocole $\mathcal{P} = (\mathsf{Setup},\mathsf{Prove},
\mathsf{Verify})$ et adversaire $\mathcal{A}$~:
\begin{enumerate}
\item $\mathcal{A}$ re\c{c}oit $pp \leftarrow \mathsf{Setup}(1^\lambda)$
\item $\mathcal{A}$ choisit $m_0,m_1 \in \mathcal{M}_\lambda$
      ($|m_0|=|m_1|$) et $c \in \mathcal{C}_\lambda$
      \hfill\textit{(phase de d\'efi)}
      \cmt{L'adversaire choisit les deux messages \`a distinguer.}
\item Le challenger tire $b \xleftarrow{\$} \{0,1\}$ et calcule
      $\sigma \leftarrow \mathsf{Prove}(pp, m_b, c)$
      \cmt{Le challenger chiffre l'un des deux messages au hasard.}
\item $\mathcal{A}$ re\c{c}oit $\sigma$ et \'emet $b' \in \{0,1\}$
      \hfill\textit{(phase de r\'eponse)}
      \cmt{L'adversaire devine lequel des deux messages a \'et\'e chiffr\'e.}
\item L'exp\'erience produit 1 si $b' = b$, sinon 0
\end{enumerate}
\begin{equation}
\mathsf{Adv}^{\mathsf{IND\text{-}CONF}}_{\mathcal{A},\mathcal{P}}(\lambda)
= \left|\Pr\!\left[
\mathsf{Exp}^{\mathsf{IND\text{-}CONF}}_{\mathcal{A},\mathcal{P}}(\lambda)
= 1\right] - \frac{1}{2}\right|
\end{equation}
\cmtFormule{L'avantage est l'\'ecart \`a la devinette al\'eatoire (probabilit\'e 1/2).}
\cmt{Un avantage n\'egligeable signifie confidentialit\'e.}
\end{definition}

\begin{remark}[Pont IND-CONF $\leftrightarrow$ $\Conf$]
\label{rem:jeu_ind_conf_link}
Pour tout distingueur $\mathcal{A}$, non born\'e en calcul, qui choisit
$(m_0,m_1,c)$~:
\begin{equation}
\mathsf{Adv}^{\mathsf{IND\text{-}CONF}}_{\mathcal{A},\mathcal{P}}(\lambda)
\;\leq\; \frac{1}{2}\max_{m_0,m_1,c}\TV(P_{\Sigma|m_0,c},P_{\Sigma|m_1,c})
\;=\; \frac{1 - \Conf(\mathcal{P})}{2},
\end{equation}
\cmtFormule{L'avantage IND-CONF est born\'e par la distinguabilit\'e statistique.}
avec \'egalit\'e pour le test de Neyman--Pearson (Helstrom) sur la paire la
plus distinguable. $\Conf$ est donc la version \emph{statistique} de
IND-CONF, au sens o\`u $\Conf \geq 1 - \negl(\lambda)$ implique un avantage n\'egligeable
pour tout adversaire. La r\'eciproque est fausse, car pour des adversaires PPT
l'avantage peut \^etre n\'egligeable (s\'ecurit\'e calculatoire) alors que
$\Conf$ est petit (chiffrement \`a cl\'e publique, par exemple).
\cmt{La confidentialit\'e statistique implique la confidentialit\'e calculatoire, pas l'inverse.}
\end{remark}

\paragraph{Jeu BIND-REL : fiabilit\'e game-based.}
\label{subsubsubsec:jeu_bind_rel}

\begin{definition}[Jeu BIND-REL]\label{def:jeu_bind_rel}
Pour un protocole $\mathcal{P}$ avec $\mathsf{Open}$ et adversaire
$\mathcal{A}$~:
\begin{enumerate}
\item $\mathcal{A}$ re\c{c}oit $pp \leftarrow \mathsf{Setup}(1^\lambda)$
\item $\mathcal{A}$ produit $(c, \sigma, s_0, s_1)$
      \cmt{L'adversaire forge une preuve et deux t\'emoins.}
\item Calculer $m_i = \mathsf{Open}(pp, \sigma, s_i)$ pour $i \in \{0,1\}$
\item $\mathcal{A}$ \emph{gagne} si et seulement si
  $\mathsf{Verify}(pp,\sigma,c)=1$,
  $m_0 \neq m_1$,
  $m_0 \neq \bot$ et $m_1 \neq \bot$
  \cmt{Il gagne si la preuve passe ET les deux messages sont distincts et valides.}
\end{enumerate}
\begin{equation}
\mathsf{Win}^{\mathsf{BIND\text{-}REL}}_{\mathcal{A},\mathcal{P}}(\lambda)
= \Pr\!\left[\mathcal{A} \text{ gagne dans }
\mathsf{Exp}^{\mathsf{BIND\text{-}REL}}_{\mathcal{A},\mathcal{P}}(\lambda)
\right]
\end{equation}
\cmtFormule{Probabilit\'e que l'adversaire produise une collision d'engagement.}
\end{definition}

\begin{remark}[Pont BIND-REL $\leftrightarrow$ $\Relver$]
\label{rem:jeu_bind_rel_link}
BIND-REL contr\^ole la \emph{solidit\'e} de l'ouverture, donc la fiabilit\'e
de v\'erification (Remarque~\ref{rem:rel_ver}). Pour un protocole complet
dont la probabilit\'e de gain BIND-REL est au plus $\varepsilon_{\mathsf{bind}}$,
une preuve accept\'ee s'ouvre sur un unique message sauf avec probabilit\'e
$\varepsilon_{\mathsf{bind}}$, de sorte que
\begin{equation}
\Relver \geq 1 - \varepsilon_{\mathsf{bind}} - \negl(\lambda).
\end{equation}
\cmtFormule{Fiabilit\'e de v\'erification born\'ee inf\'erieurement par $1 - \varepsilon_{\mathsf{bind}}$.}
BIND-REL ne dit rien de $\Relrec$, puisqu'un engagement parfaitement masquant et
calculatoirement liant a $\Relrec = 2^{-\lambda}$.
\cmt{Lier une preuve \`a un message ne revient pas \`a le r\'ev\'eler.}
\end{remark}

\paragraph{Jeu EXT-OPP : opposabilit\'e game-based.}
\label{subsubsubsec:jeu_ext_opp}

\begin{definition}[Jeu EXT-OPP]\label{def:jeu_ext_opp}
Pour un protocole $\mathcal{P}$, un interpr\`ete $\mathsf{Interpret}$
de capacit\'e $\Delta$ et une distribution $\mathcal{D}_C$ fix\'ee \`a l'avance~:
\begin{enumerate}
\item Le challenger tire $c \leftarrow \mathcal{D}_C$,
      $m \leftarrow \mathcal{M}_\lambda$
      \cmt{Le contexte est tir\'e selon la distribution du domaine.}
\item Le challenger calcule $\sigma \leftarrow \mathsf{Prove}(pp,m,c)$
\item L'interpr\`ete re\c{c}oit $(\sigma,c)$ et produit
      $v \leftarrow \mathsf{Interpret}(pp,\sigma,c,\Delta)$
      \cmt{L'interpr\`ete rend un verdict \`a partir de la preuve et du contexte.}
\item Soit $v^*$ le verdict de r\'ef\'erence, r\'ealisation de la variable $V$
      (qui peut d\'ependre du message, du contexte et des ressources
      institutionnelles, D\'efinition~\ref{def:interp_admissible})
      \cmt{Le verdict de r\'ef\'erence est celui que l'institution rendrait en pleine connaissance de cause.}
\item L'interpr\`ete \emph{gagne} si $v = v^*$
\end{enumerate}
\begin{equation}
\mathsf{Err}^{\mathsf{EXT\text{-}OPP}}_{\mathsf{Interpret},\mathcal{P},\mathcal{D}_C}(\lambda)
= \Pr\!\left[
\mathsf{Interpret} \text{ perd dans }
\mathsf{Exp}^{\mathsf{EXT\text{-}OPP}}\right]
\end{equation}
\cmtFormule{Probabilit\'e que l'interpr\`ete rende un verdict incorrect.}
\end{definition}

\begin{remark}[Pont EXT-OPP $\leftrightarrow$ $\Opp$]
\label{rem:jeu_ext_opp_link}
Par l'in\'egalit\'e de Fano appliqu\'ee \`a l'estimation de $V$ \`a partir de
$(\Sigma,C)$, $H(V|\Sigma,C) \leq h(\varepsilon_{\mathsf{ext}}) +
\varepsilon_{\mathsf{ext}}\log_2(|\mathcal{V}_\lambda|-1)$, d'o\`u
\begin{equation}
\Opp = 1 - \frac{H(V|\Sigma,C)}{H(V|C)}
\geq 1 - \frac{h(\varepsilon_{\mathsf{ext}})
  + \varepsilon_{\mathsf{ext}}\log_2(|\mathcal{V}_\lambda|-1)}
  {H(V|C)}
\label{eq:jeu_fano_opp}
\end{equation}
\cmtFormule{L'opposabilit\'e est born\'ee inf\'erieurement par $1 - \varepsilon_{\mathsf{ext}}$.}
o\`u $\varepsilon_{\mathsf{ext}} = \mathsf{Err}^{\mathsf{EXT\text{-}OPP}}$
et $h(\cdot)$ est l'entropie binaire.
Un interpr\`ete avec faible erreur d'extraction implique une forte
opposabilit\'e.
\cmt{Moins l'interpr\`ete se trompe, plus l'opposabilit\'e est grande.}
\end{remark}

\end{sloppypar}

% -----------------------------------------------------------------------
\subsection{Extension aux m\'etriques quantiques}
\label{subsec:extension_postquantique}
% -----------------------------------------------------------------------
\begin{sloppypar}
La migration vers les standards post-quantiques NIST (FIPS~203, 204, 205)~\cite{NISTPQC,NISTPostQuantum2022}
requiert une extension du cadre CLO. Cette section esquisse les modifications
n\'ecessaires, d\'evelopp\'ees en d\'etail dans la section~\ref{sec:cro_pq}.

\subsubsection{Confidentialit\'e quantique}
\label{subsubsec:conf_quantique}

En pr\'esence d'adversaires quantiques, la d\'efinition de confidentialit\'e
doit consid\'erer les mesures POVM sur les \'etats quantiques de
preuve~\cite{Watrous2009,Boneh2011,Zhandry}. Comme dans le cas classique, la
m\'etrique retenue est \emph{statistique}, puisqu'elle quantifie la meilleure
distinction possible, sans restriction de calcul.

\begin{definition}[Confidentialit\'e quantique]\label{def:conf_quantique_clo}
Pour le protocole $\mathcal{P}$ avec param\`etre de s\'ecurit\'e $\lambda$, soit
$\rho_\Sigma^{m,c}$ l'\'etat quantique de preuve pour le message $m$ dans le
contexte $c$. La \emph{confidentialit\'e quantique} est
\begin{equation}
\Conf^q(\mathcal{P}) := 1 - \max_{m_0,m_1 \in \mathcal{M}_\lambda,\; c \in \mathcal{C}_\lambda}
\tfrac{1}{2}\bigl\|\rho_\Sigma^{m_0,c} - \rho_\Sigma^{m_1,c}\bigr\|_1 .
\label{eq:conf_q}
\end{equation}
Par le th\'eor\`eme de Helstrom, $\tfrac12\|\rho_0-\rho_1\|_1$ est l'avantage
maximal de toute mesure POVM pour distinguer $\rho_0$ de $\rho_1$.
\end{definition}

\begin{remark}[Compatibilit\'e avec la confidentialit\'e classique]
\label{rem:compatibilite_conf}
Pour les protocoles classiques o\`u $\Sigma$ prend des valeurs dans un alphabet
fini $\mathcal{S}$ (i.e., $\rho_\Sigma^{m,c} = \sum_{s} P_{\Sigma|M,C}(s|m,c)\,|s\rangle\langle s|$
est diagonal), la distance de trace se r\'eduit \`a la distance de variation totale~\cite{Levin,Pinsker1964}~:
\begin{equation}
\tfrac{1}{2}\|\rho_\Sigma^{m_0,c} - \rho_\Sigma^{m_1,c}\|_1
= \mathrm{TV}\!\bigl(P_{\Sigma|m_0,c},\, P_{\Sigma|m_1,c}\bigr)
\end{equation}
Donc $\Conf^q$ se r\'eduit \`a $\Conf$ du cadre classique~\cite{MINKA2026CLO}, ce qui fait de la
d\'efinition une g\'en\'eralisation stricte. Les notions calculatoires
(adversaires QPT) rel\`event des jeux de s\'ecurit\'e, comme IND-CONF dans le
cas classique (Remarque~\ref{rem:jeu_ind_conf_link}).
\end{remark}

\subsubsection{Perte d'interpr\'etabilit\'e quantique}
\label{subsubsec:perte_interpretabilite}

\begin{definition}[Perte d'interpr\'etabilit\'e quantique]\label{def:perte_interp_quantique}
Soit $\Lambda$ une application CPTP repr\'esentant les effets d'un canal
quantique sur le support de la preuve (stockage, transmission,
v\'erification). La perte d'interpr\'etabilit\'e quantique $\eta_q \in [0,1]$
est le plus petit r\'eel tel que
\begin{equation}
I(V;\Lambda(\Sigma)|C) \;\geq\; I(V;\Sigma|C) - \eta_q\, H(C).
\label{eq:def_eta_q}
\end{equation}
\end{definition}

L'opposabilit\'e apr\`es passage dans le canal est alors
\begin{equation}
\Opp^q(\mathcal{P}) = \frac{I(V;\Lambda(\Sigma)|C)}{H(V|C)} \;\geq\; \Opp(\mathcal{P}) - \eta_q\,\frac{H(C)}{H(V|C)},
\label{eq:opp_q}
\end{equation}
une minoration qui quantifie la d\'egradation s\'emantique dans les canaux
quantiques~\cite{Holevo1973,Wilde2017}. Par l'in\'egalit\'e de traitement des
donn\'ees, $I(V;\Lambda(\Sigma)|C) \leq I(V;\Sigma|C)$. Un canal ne peut donc que
r\'eduire l'opposabilit\'e, tandis que $\eta_q$ mesure l'ampleur maximale de cette
r\'eduction, rapport\'ee \`a la richesse du contexte.

\begin{remark}[Interpr\'etation physique]
\label{rem:interpretation_physique}
$\eta_q$ quantifie la quantit\'e de contenu s\'emantique perdue lorsque la
preuve cryptographique passe par des canaux quantiques~\cite{Zurek2003,Schumacher1997}. Pour les syst\`emes
dont la preuve est une cha\^ine de bits classique, y compris les signatures
post-quantiques standardis\'ees, $\Lambda$ est l'identit\'e et $\eta_q = 0$. Pour
un support ou un v\'erificateur quantique, $\eta_q > 0$ en raison du bruit
(d\'ecoh\'erence, erreurs de portes)~\cite{Holevo1973,Preskill2018}.
\end{remark}

Ces extensions seront formalis\'ees et exploit\'ees \`a la section~\ref{sec:cro_pq} pour \'etablir le Trilemme CRO conditionnel quantique.
\end{sloppypar}

% -----------------------------------------------------------------------
% Synth\`ese
% -----------------------------------------------------------------------
\begin{boiteArbitre}
L'arbitre referme le dossier CLO et ach\`eve ainsi l'\'etape \textbf{F}. Cette section lui a donn\'e un
langage pour mesurer ce qu'il vit, qu'aucun corpus pr\'ec\'edent ne lui avait fourni.

Il retient trois contributions, dans l'ordre o\`u elles se sont
construites. Le mod\`ele de protocole d'abord (deux couches et deux questions
distinctes, portant sur la validit\'e de la preuve puis sur ce qu'elle signifie pour lui).
Les trois m\'etriques ensuite ($\Conf$, $\Rel$, $\Opp$), dont la derni\`ere
le concerne directement~; il sait d\'esormais qu'elle
n'a de sens que si son propre jugement ne se r\'eduit pas \`a un calcul
m\'ecanique sur la preuve. Le goulot d'\'etranglement contextuel enfin,
qui lui a appris que le plafond est fix\'e par la finesse de ses verdicts.

La section~\ref{sec:cro_classique} l'attend. Il va falloir maintenant
savoir ce qui, sous des hypoth\`eses explicites sur l'institution, borne
le produit des trois m\'etriques.
\end{boiteArbitre}